\clearpage\thispagestyle{empty}
\pdfbookmark[0]{Quellen zur Tafel für die Sonnenkoordinaten (1812)}{sonnen1812-sources}
\begin{center}\textbf{Quellen zur Tafel für die Sonnenkoordinaten (1812)}\end{center}
\noindent Diese KI-Übertragung in heutiges Deutsch folgt Carl Friedrich Gauß’ deutschem Text in den \href{https://gdz.sub.uni-goettingen.de/id/PPN236007467}{\textit{Werken}, Band~VI} (Göttingen, 1874), S.~172--180. Sie enthält das Rechenbeispiel und die Tafeln der Epochen, der Sonnenkoordinaten und ihrer säkularen Änderungen.
\par\smallskip\noindent Der \href{https://www.digitale-sammlungen.de/en/view/bsb10538614?page=29}{Erstdruck} erschien als Beitrag~II, \textit{Über die Tafel für die Sonnen-Coordinaten in Beziehung auf den Äquator}, im Januarheft 1812 der von Franz Xaver von Zach herausgegebenen \textit{Monatlichen Correspondenz zur Beförderung der Erd- und Himmels-Kunde}, Band~25, S.~23--36. Die \href{https://www.digitale-sammlungen.de/en/view/bsb10538614?page=34}{Tafeln auf S.~28--36} nennen \href{https://www.deutsche-biographie.de/gnd134127897.html#ndbcontent}{Friedrich Ludwig Wachter} als Bearbeiter der Berechnungen. Gauß verwendet im Beispiel den Zeitpunkt einer Beobachtung von Zachs und gibt die unabhängig von Lindenau berechneten Werte an.
\par\smallskip\noindent Übertragung und Satz wurden mit KI erstellt.
\clearpage
